% Part: computability
% Chapter: computability-theory
% Section: russells-paradox

\documentclass[../../../include/open-logic-section]{subfiles}

\begin{document}

\olfileid{cmp}{thy}{rus}
\olsection{مقایسه با پارادوکس راسل}

مقایسهٔ استدلال‌های این بخش با پارادوکس راسل و توجه به تفاوت آن‌ها
آموزنده است:
\begin{enumerate}
\item پارادوکس راسل: بگذارید $S=\Setabs{x}{x\notin x}$. آنگاه
  $S\in S$ اگر و تنها اگر $S\notin S$؛ تناقض.

  \emph{نتیجه:} چنین مجموعهٔ~$S$ای وجود ندارد. فرض وجود «مجموعهٔ همهٔ
  مجموعه‌ها» با دیگر اصول موضوع نظریهٔ مجموعه‌ها ناسازگار است.

\item تغییری در پارادوکس راسل: بگذارید $F$ «تابعی» از مجموعهٔ همهٔ
  توابع به $\{0,1\}$ باشد که چنین تعریف شده است:
  \[
  F(f) =
  \begin{cases}
    1 & \text{اگر $f$ در دامنهٔ $f$ باشد و $f(f) = 0$} \\
    0 & \text{در غیر این صورت.}
  \end{cases}
  \]
  استدلالی همانند نشان می‌دهد $F(F)=0$ اگر و تنها اگر $F(F)=1$؛
  تناقض.

  \emph{نتیجه:} $F$ تابع نیست. «مجموعهٔ همهٔ توابع» بزرگ‌تر از آن است
  که دامنهٔ یک تابع باشد.

\item استدلال قطری‌سازی: بگذارید $f_0$، $f_1$، \dots برشماری توابع
  محاسبه‌پذیر جزئی باشد و $G\colon\Nat\to\{0,1\}$ را چنین تعریف کنید:
  \[
  G(x) =
  \begin{cases}
    1 & \text{اگر $f_x(x)\downarrow = 0$} \\
    0 & \text{در غیر این صورت.}
  \end{cases}
  \]
  اگر $G$ محاسبه‌پذیر باشد، برای یک $k$ همان تابع $f_k$ است. اما در
  این صورت $G(k)=1$ اگر و تنها اگر $G(k)=0$؛ تناقض.

  \emph{نتیجه:} $G$ محاسبه‌پذیر نیست. توجه کنید که بنا بر اصول موضوع
  نظریهٔ مجموعه‌ها، $G$ همچنان یک تابع است؛ اینجا پارادوکسی در کار
  نیست، بلکه تنها ابهامی رفع شده است.
\end{enumerate}

سخن‌گفتن از توابع جزئی، توابع محاسبه‌پذیر، توابع محاسبه‌پذیر جزئی و
جز آن ممکن است گیج‌کننده باشد. مجموعهٔ همهٔ توابع جزئی از $\Nat$ به
$\Nat$ مجموعه‌ای بزرگ از اشیاست. برخی تام‌اند، برخی محاسبه‌پذیرند،
برخی هم تام و هم محاسبه‌پذیرند و برخی هیچ‌یک نیستند. به یاد داشته
باشید که وقتی بدون قید از «تابع» سخن می‌گوییم، مقصود تابعی تام است.
پس داریم:
\begin{enumerate}
\item توابع محاسبه‌پذیر؛
\item توابع محاسبه‌پذیر جزئی که تام نیستند؛
\item توابعی که محاسبه‌پذیر نیستند؛
\item توابع جزئی‌ای که نه تام‌اند و نه محاسبه‌پذیر.
\end{enumerate}
برای مرتب‌کردن این رده‌ها شاید کمک کند مربعی بزرگ بکشید که نمایندهٔ
همهٔ توابع جزئی از $\Nat$ به $\Nat$ باشد و سپس دو ناحیهٔ هم‌پوشان
را به‌ترتیب برای توابع تام و توابع جزئیِ محاسبه‌پذیر مشخص کنید.
تمرین خوبی است که ببینید آیا می‌توانید در هر یک از نواحی حاصل در
نمودار، شیئی توصیف کنید.

\end{document}
